\section{Problem and Security Model}
\label{sec:threat-model}

\paragraph{System and enforcement boundary.}
A user gives an agent an authenticated task $q$.  The agent may read messages,
documents, webpages, or database records before proposing a structured tool
call $c=(t,\bar{x})$.  The LLM is an untrusted control-plane component: its plan,
effect description, and call arguments are proposals rather than authority.
Executing the fully instantiated call from pre-state $s$ may commit several
externally visible effects, such as creating an event, inviting principals, or
changing its visibility.  The enforcement point lies after argument
totalization but immediately before those effects commit.  The executor may
commit only the exact call associated with the monitor's decision.

\paragraph{Threat model.}
The adversary controls content returned by one or more untrusted reads.  It may
embed instructions or attacker-chosen values, induce the agent to add a
resource, target, payload, operation, or effect, and carry this influence across
a multi-step trajectory.  We assume the agent may follow the injected
instruction completely; security does not rely on model alignment.  The
adversary cannot modify the authenticated task, forge trusted provenance,
change the registered descriptor or policy, compromise the monitor, or replace
a checked call before execution.  It also cannot modify the tool implementation
after registration.  The trusted harness marks untrusted-content boundaries;
the AgentDojo adapter derives typed provenance from those boundaries.

\paragraph{Security objective.}
At the architectural level, let $u=(t,\bar{x},s,h)$ denote an execution context
containing a tool, totalized arguments, pre-state, and typed provenance evidence.
Let $\mu(u)$ denote its concrete committed effects.  A policy context $P$
may additionally depend on trusted execution metadata $\gamma(u)$, including
effect provenance.  Writing $\omega(u)=(\mu(u),\gamma(u))$, it induces an ideal
per-commit predicate $I_P(\omega(u))$.  The desired commit condition is
$I_P(\omega(u))=1$; multiset containment $\mu(u)\preceq B$ is one
concrete instance.  This paper addresses the representation obligation that
the monitor's view must preserve every distinction needed to evaluate that
predicate.  The authenticated user, harness, or an external policy component
supplies $P$; the planner may propose values but cannot create or enlarge the
authority against which those values are checked.

Our runtime prototype instantiates one bounded policy, \emph{provenance-origin
confinement}.  An untrusted observation may provide data, but it may not be the
sole origin of a registered effect or effect-bearing value absent from the
authenticated task.  Thus the evaluated monitor tests whether a proposed call crosses this source
boundary; it does not establish that the user is globally entitled to the
action.  ACLs, delegation, quotas, and organization policy remain independent
authority predicates that could consume the same effect representation.

\paragraph{Trust assumptions and scope.}
Our conditional security argument requires a descriptor that covers in-scope
effects, unforgeable provenance, sound task grounding, and a correct authority
predicate.  Every effectful call must be mediated, and the checked call and
policy snapshot must govern the commit.  Registration additionally assumes
that the sandbox oracle observes the effects exercised by its interventions
and that runtime invokes the registered implementation.
If the planner can alter $P$, relabel provenance, or replace the checked call,
the stated confinement result does not hold.
We do not address compromise of the host, monitor, provenance adapter, or tool
implementation; post-registration code changes; effects outside the observed
state boundary; output-only attacks; or dangerous actions explicitly requested
by the authenticated user.
